\documentclass[11pt,a4paper]{scrartcl}

\usepackage{amssymb}

\usepackage[utf8]{inputenc}
\usepackage[T1]{fontenc}
\newcommand{\ats}[1]{\par\noindent\textit{Atsakymas:} #1}
\usepackage{cancel}
\usepackage{tikz}
\usepackage{lmodern} 
\usepackage{amsmath,amssymb,amsthm}
\usepackage{graphicx}
\usepackage[dvipsnames,svgnames]{xcolor}
\usepackage{hyperref}
\usepackage{mathrsfs}
\usepackage[shortlabels]{enumitem}
\usepackage[obeyFinal,textsize=scriptsize,shadow,loadshadowlibrary]{todonotes}
\usepackage{mathtools}
\usepackage{microtype}

%%% Paragraph layout
\setlength{\parindent}{0pt}
\setlength{\parskip}{0.6\baselineskip}

%%% Colored sections
\renewcommand*{\sectionformat}%
  {\color{purple}\S\thesection\enskip}
\renewcommand*{\subsectionformat}%
  {\color{purple}\S\thesubsection\enskip}
\renewcommand*{\subsubsectionformat}%
  {\color{purple}\S\thesubsubsection\enskip}
\addtokomafont{paragraph}{\color{orange!35!black}\P\ }
\KOMAoptions{numbers=noenddot}

%%% Title
\addtokomafont{subtitle}{\Large}
\setkomafont{author}{\Large\scshape}
\setkomafont{date}{\Large\normalsize}
% Neigiamas vertikalus tarpas, kuris pakelia dokumento antraštę į viršų. 
% Galite keisti -2.5cm į -3cm (aukščiau) arba -1.5cm (žemiau).
\titlehead{\vspace*{-7.5cm}} 

% PRIDĖTA: Automatiškai perrašome datą į mokyklos pavadinimą
\AtBeginDocument{%
  \date{\normalsize Vilniaus licėjus}
}

%%% Page setup
\usepackage[headsepline]{scrlayer-scrpage}
\renewcommand{\headfont}{}
\addtolength{\textheight}{3.14cm}
\setlength{\footskip}{0.5in}
\setlength{\headsep}{10pt}
% Puslapio antraštė turi rodyti tik konkrečius dokumento duomenis.
% Nustatome juos aiškiai, kad neatsirastų "\@author" / "\@title" arba senų reikšmių.
\makeatletter
\AtBeginDocument{%
  \ihead{\footnotesize\textbf{Andrius Berniukevičius}}%
  \ohead{\footnotesize\textbf{\@title}}%
}
\makeatother
\automark{section}
\chead{}
\cfoot{\pagemark}

%%% Macros
\providecommand{\ol}{\overline}
\providecommand{\ul}{\underline}
\providecommand{\wt}{\widetilde}
\providecommand{\wh}{\widehat}
\providecommand{\eps}{\varepsilon}
\providecommand{\half}{\frac{1}{2}}
\providecommand{\inv}{^{-1}}
\newcommand{\dang}{\measuredangle} %% Directed angle
\newcommand{\vocab}[1]{\textbf{\color{blue}\sffamily #1}}
\providecommand{\CC}{\mathbb C}
\providecommand{\FF}{\mathbb F}
\providecommand{\NN}{\mathbb N}
\providecommand{\QQ}{\mathbb Q}
\providecommand{\RR}{\mathbb R}
\providecommand{\ZZ}{\mathbb Z}
\providecommand{\ts}{\textsuperscript}
\providecommand{\dg}{^\circ}
\providecommand{\ii}{\item}
\providecommand{\defeq}{\coloneqq}
\DeclareMathOperator*{\lcm}{lcm}
\DeclareMathOperator*{\argmin}{arg min}
\DeclareMathOperator*{\argmax}{arg max}
\providecommand{\hrulebar}{\par
\hspace{\fill}\rule{0.95\linewidth}{.7pt}\hspace{\fill}
\par\nointerlineskip \vspace{\baselineskip}}

%%% Optional colored theorems
\usepackage{thmtools}
\usepackage[framemethod=TikZ]{mdframed}
\mdfdefinestyle{mdbluebox}{%
  roundcorner=10pt,
  linewidth=1pt,
  skipabove=12pt,
  innerbottommargin=9pt,
  skipbelow=2pt,
  linecolor=blue,
  nobreak=true,
  backgroundcolor=TealBlue!5,
}
\declaretheoremstyle[
  headfont=\sffamily\bfseries\color{MidnightBlue},
  mdframed={style=mdbluebox},
  headpunct={\\[3pt]},
  postheadspace={0pt}
]{thmbluebox}
\mdfdefinestyle{mdredbox}{%
  linewidth=0.5pt,
  skipabove=12pt,
  frametitleaboveskip=5pt,
  frametitlebelowskip=0pt,
  skipbelow=2pt,
  frametitlefont=\bfseries,
  innertopmargin=4pt,
  innerbottommargin=8pt,
  nobreak=true,
  backgroundcolor=Salmon!5,
  linecolor=RawSienna,
}
\declaretheoremstyle[
  headfont=\bfseries\color{RawSienna},
  mdframed={style=mdredbox},
  headpunct={\\[3pt]},
  postheadspace={0pt},
]{thmredbox}
\mdfdefinestyle{mdgreenbox}{%
  skipabove=8pt,
  linewidth=2pt,
  rightline=false,
  leftline=true,
  topline=false,
  bottomline=false,
  linecolor=ForestGreen,
  backgroundcolor=ForestGreen!5,
}
\declaretheoremstyle[
  headfont=\bfseries\sffamily\color{ForestGreen!70!black},
  bodyfont=\normalfont,
  spaceabove=2pt,
  spacebelow=1pt,
  mdframed={style=mdgreenbox},
  headpunct={ --- },
]{thmgreenbox}
\mdfdefinestyle{mdblackbox}{%
  skipabove=8pt,
  linewidth=3pt,
  rightline=false,
  leftline=true,
  topline=false,
  bottomline=false,
  linecolor=black,
  backgroundcolor=RedViolet!5!gray!5,
}
\declaretheoremstyle[
  headfont=\bfseries,
  bodyfont=\normalfont\small,
  spaceabove=0pt,
  spacebelow=0pt,
  mdframed={style=mdblackbox}
]{thmblackbox}

%% Aplinkos su rėmeliais (thmtools)
\declaretheorem[style=thmbluebox,name=Teorema]{theorem}
\declaretheorem[style=thmbluebox,name=Lema]{lemma}
\declaretheorem[style=thmbluebox,name=Savybė]{property}
\declaretheorem[style=thmbluebox,name=Teiginys]{proposition}
\declaretheorem[style=thmbluebox,name=Išvada]{corollary}
\declaretheorem[style=thmbluebox,name=Teorema,numbered=no]{theorem*}
\declaretheorem[style=thmbluebox,name=Lema,numbered=no]{lemma*}
\declaretheorem[style=thmbluebox,name=Teiginys,numbered=no]{proposition*}
\declaretheorem[style=thmbluebox,name=Išvada,numbered=no]{corollary*}
\declaretheorem[style=thmgreenbox,name=Algoritmas]{algorithm}
\declaretheorem[style=thmgreenbox,name=Algoritmas,numbered=no]{algorithm*}
\declaretheorem[style=thmgreenbox,name=Teiginys]{claim}
\declaretheorem[style=thmgreenbox,name=Teiginys,numbered=no]{claim*}
\declaretheorem[style=thmredbox,name=Pavyzdys]{example}
\declaretheorem[style=thmredbox,name=Pavyzdys,numbered=no]{example*}
\declaretheorem[style=thmblackbox,name=Pastaba]{remark}
\declaretheorem[style=thmblackbox,name=Pastaba,numbered=no]{remark*}

%% Standartinės amsthm aplinkos (Apibrėžimai ir užduotys)
\theoremstyle{definition}
\ifcsname conjecture\endcsname\else\newtheorem{conjecture}{Hipotezė}\fi
\ifcsname definition\endcsname\else\newtheorem{definition}{Apibrėžimas}\fi
\ifcsname fact\endcsname\else\newtheorem{fact}{Faktas}\fi
\ifcsname answer\endcsname\else\newtheorem{answer}{Atsakymas}\fi
\ifcsname ques\endcsname\else\newtheorem{ques}{Klausimas}\fi
\ifcsname exercise\endcsname\else\newtheorem{exercise}{Užduotis}\fi
\ifcsname problem\endcsname\else\newtheorem{problem}{Uždavinys}[section]\fi
\ifcsname conjecture*\endcsname\else\newtheorem*{conjecture*}{Hipotezė}\fi
\ifcsname definition*\endcsname\else\newtheorem*{definition*}{Apibrėžimas}\fi
\ifcsname fact*\endcsname\else\newtheorem*{fact*}{Faktas}\fi
\ifcsname answer*\endcsname\else\newtheorem*{answer*}{Atsakymas}\fi
\ifcsname ques*\endcsname\else\newtheorem*{ques*}{Klausimas}\fi
\ifcsname exercise*\endcsname\else\newtheorem*{exercise*}{Užduotis}\fi
\ifcsname problem*\endcsname\else\newtheorem*{problem*}{Uždavinys}\fi

\newcounter{answerssection}
\newcommand{\answerssection}{%
  \setcounter{answerssection}{\value{section}}%
  \section{Atsakymai}%
  \setcounter{problem}{0}%
  \renewcommand{\theproblem}{\arabic{answerssection}.\arabic{problem}}%
}

%% GALUTINIS SPRENDIMAS PROOF APLINKAI
%% --- PRADŽIA: Įrodymo aplinkos sutvarkymas ---
\makeatletter

% 1. Užtikriname, kad žodis būtų "Įrodymas", net jei "babel" paketas bando jį perrašyti
\AtBeginDocument{%
  \renewcommand{\proofname}{Įrodymas}%
  \@ifpackageloaded{babel}{%
    \addto\captionslithuanian{\renewcommand{\proofname}{Įrodymas}}%
  }{}%
}

% 2. Priverstinai pašaliname scrartcl klasės bold šriftą ir paliekame TIK italic
% 3. Sujungiame su tuščiu kvadratėliu dešiniajame kampe.
\renewcommand{\qedsymbol}{\ensuremath{\Box}}
\renewenvironment{proof}[1][\proofname]{\par
  \pushQED{\qed}%
  \normalfont \topsep6\p@\@plus6\p@\relax
  \trivlist
  \item[\hskip\labelsep
        \normalfont\itshape % Priverčia naudoti tik pasvirą šriftą, kaip nuotraukoje
    #1\@addpunct{.}]\ignorespaces
}{%
  \popQED\endtrivlist\@endpefalse
}
\newenvironment{solution}{\par
  \normalfont \topsep6\p@\@plus6\p@\relax
  \trivlist
  \item[\hskip\labelsep
        \normalfont\itshape
    \textit{Sprendimas}\@addpunct{.}]\ignorespaces
}{%
  \endtrivlist\@endpefalse
}
\newcommand{\ans}[1]{\par\noindent\textit{Atsakymas:} #1}
\makeatother


\title{Kaip mąsto olimpiadininkas?}
\author{Andrius Berniukevičius}

\newenvironment{solution}
    {\begin{list}{}{\setlength{\leftmargin}{10pt}\setlength{\rightmargin}{10pt}}\item[]}
    {\end{list}}

\begin{document}

\maketitle

\section{Svarbiausia mintis}

Sprendžiant olimpiadinį uždavinį vyksta du skirtingi dalykai:

\[
\boxed{\text{sprendimo paieška}}
\qquad\text{ir}\qquad
\boxed{\text{sprendimo įrodymas}}.
\]

Ieškant sprendimo visiškai normalu:

\begin{itemize}
    \item bandyti mažus atvejus;
    \item skaičiuoti konkrečius pavyzdžius;
    \item piešti brėžinius;
    \item spėti dėsningumą;
    \item bandyti kelias skirtingas idėjas;
    \item klysti ir pradėti iš naujo.
\end{itemize}

Tačiau galutinis sprendimas turi atsakyti į kitą klausimą:

\[
\text{„Kodėl tai \textbf{būtinai} yra tiesa?“}
\]

Todėl labai svarbu atskirti:

\[
\text{„Atrodo, kad taip yra“}
\]
nuo
\[
\text{„Taip turi būti, nes \ldots“}.
\]

Galutinis užrašytas sprendimas gali atrodyti visai kitaip negu kelias,
kuriuo jį atradote. Tai yra normalu. Kad tai pamatytumėte, palyginkite, kaip atrodo tas pats uždavinys juodraštyje ir švarraštyje.

\bigskip
\textbf{Pavyzdys:} \textit{Įrodykite, kad su visais sveikaisiais $n$, skaičius $n^2+n$ yra lyginis.}

\noindent
\textbf{Juodraštyje (sprendimo paieška):}
\begin{itemize}
    \item Jei $n=1$, tai $1^2+1 = 2$ (lyginis).
    \item Jei $n=2$, tai $2^2+2 = 6$ (lyginis).
    \item Jei $n=3$, tai $3^2+3 = 12$ (lyginis).
    \item Aha! Jei $n$ nelyginis, $n^2$ nelyginis, o nelyginis + nelyginis = lyginis.
    \item Jei $n$ lyginis, tai lyginis + lyginis = lyginis. Viskas aišku!
\end{itemize}

\noindent
\textbf{Švarraštyje (olimpiados darbe):}
\begin{quote}
Išskaidome reiškinį dauginamaisiais: $n^2+n = n(n+1)$. 
Kadangi $n$ ir $n+1$ yra du iš eilės einantys sveikieji skaičiai, lygiai vienas iš jų privalo būti lyginis. Lyginio ir bet kokio sveikojo skaičiaus sandauga visada yra lyginis skaičius, todėl $n(n+1)$ yra lyginis visiems sveikiesiems $n$. $\square$
\end{quote}

\section{Kaip pradėti galvoti apie olimpiadinį uždavinį?}

Įprastame mokykliniame uždavinyje dažnai iš anksto aišku, kokį metodą
reikia taikyti. Olimpiadiniame uždavinyje svarbi sprendimo dalis yra
\textbf{pačiam atpažinti idėją}. Todėl pirmasis klausimas neturėtų būti

\[
\text{„Kokią formulę čia reikia panaudoti?“}
\]
Geriau klausti:
\[
\text{„Ką galiu išvesti iš to, kas duota?“}
\]
Sprendžiant verta sau užduoti tokius klausimus:

\begin{itemize}
    \item Kas tiksliai yra duota?
    \item Ką tiksliai reikia įrodyti arba rasti?
    \item Ar galiu sąlygą užrašyti kita forma?
    \item Kas nutinka mažais arba paprastais atvejais?
    \item Ar objektus galima suskirstyti į kelis atvejus?
    \item Kas būtų, jeigu teiginys būtų neteisingas?
    \item Ko pakaktų, kad gaučiau tai, ką reikia įrodyti?
\end{itemize}

\section{Eksperimentas padeda atrasti, bet neįrodo}

Mažų atvejų tikrinimas yra labai geras įprotis. Jis gali padėti pastebėti
dėsningumą arba atspėti atsakymą. Tačiau keli pavyzdžiai dar nėra įrodymas. Pavyzdžiui, reiškinys

\[
n^2+n+41
\]
yra pirminis, kai
\[
n=0,1,2,\ldots,39.
\]
Po tiek daug pavyzdžių labai lengva patikėti, kad jis visada pirminis.
Tačiau, kai $n=40$,
\[
40^2+40+41=1681=41^2,
\]
todėl gautas skaičius jau nėra pirminis. Taigi eksperimentavimas turi labai svarbią paskirtį:

\[
\boxed{\text{pavyzdžiai padeda atrasti, ką verta bandyti įrodyti.}}
\]
Bet pats įrodymas turi paaiškinti, \textbf{kodėl} teiginys galioja visais
reikalingais atvejais.

\section{Pavyzdys nėra įrodymas}

Tarkime, norime parodyti, kad dviejų nelyginių sveikųjų skaičių suma yra
lyginė. Galime patikrinti:

\[
3+5=8,\qquad 7+11=18,\qquad 101+203=304.
\]

Visais trimis atvejais teiginys teisingas. Tačiau sveikųjų skaičių yra
be galo daug, todėl šie pavyzdžiai dar neįrodo bendro teiginio. Norėdami įrodyti jį visiems nelyginiams skaičiams, naudojame bendrą
nelyginio sveikojo skaičiaus formą
\[
2k+1,
\]
kur $k\in\mathbb Z$.

\begin{example}
Įrodykite, kad dviejų nelyginių sveikųjų skaičių suma yra lyginė.
\end{example}

\textit{Prieš skaitydami sprendimą, pabandykite jį užrašyti patys.}

\begin{solution}
\textbf{Sprendimas}

Tegul
\[
a=2m+1,\qquad b=2n+1,
\]
kur $m,n\in\mathbb Z$. Tuomet
\[
a+b=(2m+1)+(2n+1)=2(m+n+1).
\]
Kadangi $m+n+1$ yra sveikasis skaičius, tai $a+b$ dalijasi iš $2$. Vadinasi, $a+b$ yra lyginis. \hfill $\square$
\end{solution}

Svarbiausia čia ne pats skaičiavimas. Svarbiausia buvo pereiti nuo kelių
konkrečių pavyzdžių prie \textbf{bendro argumento}.

\section{Teiginio kryptis}

Daugelis matematinių teiginių turi formą
\[
A\Longrightarrow B.
\]
Tai reiškia: jeigu $A$ yra tiesa, tai turi būti tiesa ir $B$. Labai svarbu nepamiršti, kad iš
\[
A\Longrightarrow B
\]
nebūtinai seka
\[
B\Longrightarrow A.
\]
Pavyzdžiui,
\[
6\mid n\Longrightarrow 3\mid n
\]
yra teisingas teiginys. Tačiau atvirkštinis teiginys
\[
3\mid n\Longrightarrow 6\mid n
\]
yra klaidingas: užtenka paimti $n=9$.

Olimpiadoje galima parašyti daug teisingų formulių ir vis tiek neįrodyti
to, ko buvo prašoma, jeigu įrodoma ne ta implikacijos kryptis.

\section{Kontrapozicija}

Kartais teiginį
\[
A\Longrightarrow B
\]
patogiau įrodyti parodant lygiavertį teiginį
\[
\neg B\Longrightarrow \neg A.
\]
Tai vadinama \textbf{kontrapozicija}.
Kodėl tai veikia? Nes teiginys „jeigu $A$, tai $B$“ yra tas pats, kas
„jeigu $B$ neįvyksta, tai $A$ neįvyksta“. Kitaip tariant, jei kažkas
nesilaiko $B$, tai automatiškai nereiškia, kad $A$ buvo teisinga. Pavyzdžiui, paprastas sakinys:

\[
\text{Jeigu lietus lyja, tai gatvė yra šlapia.}
\]
yra ekvivalentus su:
\[
\text{Jeigu gatvė nėra šlapia, tai lietus nelyja.}
\]
Abu sakiniai reiškia tą patį. Todėl kartais lengviau įrodyti
neigiamą versiją, nei pačią originalią formulę. Štai kodėl olimpiadoje dažnai apsuka teiginį: užuot bandę tiesiogiai
parodyti, kad iš $A$ seka $B$, jie įrodo, kad iš $\neg B$ seka $\neg A$.
Tai dažnai suveda problemą prie paprastesnių ar labiau matomų atvejų.

\begin{example}
Įrodykite: jei $n^2$ yra lyginis, tai $n$ yra lyginis.
\end{example}

\begin{solution}
\textbf{Sprendimas}

Įrodysime kontrapoziciją: jeigu $n$ yra nelyginis, tai $n^2$ yra
nelyginis. Tegul
\[
n=2k+1,
\]
kur $k\in\mathbb Z$. Tuomet
\[
n^2=(2k+1)^2=4k^2+4k+1
=2(2k^2+2k)+1.
\]
Taigi $n^2$ yra nelyginis. Vadinasi,
\[
n^2\text{ lyginis}\Longrightarrow n\text{ lyginis}. \hfill \square
\]
\end{solution}

Apie kontrapoziciją verta pagalvoti tada, kai sunku tiesiogiai įrodyti
$A\Longrightarrow B$, tačiau lengva suprasti, kas nutiktų, jeigu $B$
būtų netiesa.

\section{Įrodymas prieštara}

Dar vienas svarbus metodas yra \textbf{įrodymas prieštara}. Jo idėja:

\begin{enumerate}
    \renewcommand{\labelenumi}{\textbf{\theenumi.}}
    \item Tarkime, kad norimas teiginys yra neteisingas.
    \item Logiškai tęskime šią prielaidą.
    \item Parodykime, kad gauname neįmanomą rezultatą.
    \item Vadinasi, pradinė prielaida buvo klaidinga.
\end{enumerate}

\begin{example}
Įrodykite, kad nėra didžiausio sveikojo skaičiaus.
\end{example}

\begin{solution}
\textbf{Sprendimas}

Tarkime priešingai: egzistuoja didžiausias sveikasis skaičius $N$. Tačiau skaičius
\[
N+1
\]
taip pat yra sveikasis ir
\[
N+1>N.
\]
Tai prieštarauja prielaidai, kad $N$ yra didžiausias sveikasis skaičius.
Vadinasi, didžiausio sveikojo skaičiaus nėra. \hfill $\square$
\end{solution}

\section{Kontrpavyzdys}

Norint įrodyti teiginį, kuris turi galioti \textbf{visiems} objektams,
reikia bendro argumento. Tačiau norint parodyti, kad toks teiginys yra neteisingas, pakanka rasti
\textbf{vieną} atvejį, kuriam jis negalioja. Toks atvejis vadinamas \textbf{kontrpavyzdžiu}.

\begin{example}
Ar teisingas teiginys:
\[
\text{„Jeigu }ab\text{ yra lyginis, tai }a\text{ ir }b\text{ yra lyginiai“?}
\]
\end{example}

\begin{solution}
\textbf{Sprendimas}
Ne. Paimkime
\[
a=2,\qquad b=3.
\]
Tuomet
\[
ab=6
\]
yra lyginis, tačiau $b=3$ yra nelyginis.

Taigi radome kontrpavyzdį, todėl teiginys neteisingas. \hfill $\square$
\end{solution}

Svarbu:
\begin{itemize}
    \item keli teisingi pavyzdžiai dar neįrodo bendro teiginio;
    \item vienas kontrpavyzdys jau paneigia bendrą teiginį.
\end{itemize}

\section{„Raskite visus“ reiškia daugiau negu „raskite kelis“}

Jeigu uždavinys prašo \textbf{rasti visus} skaičius, poras ar kitus
objektus, reikia atlikti tris darbus:

\begin{enumerate}
    \renewcommand{\labelenumi}{\textbf{\theenumi.}}
    \item rasti galimus kandidatus;
    \item įrodyti, kad kitų kandidatų nėra;
    \item patikrinti, kad rasti kandidatai iš tikrųjų tenkina pradinę sąlygą.
\end{enumerate}

\begin{example}
Raskite visus sveikuosius skaičius $n$, kuriems trupmenos
\[
\frac{n^2+3n+1}{n+1}
\]
reikšmė yra sveikasis skaičius.
\end{example}

\begin{solution}
\textbf{Sprendimas}

Pirmiausia pastebime, kad $n\neq -1$. Pertvarkome skaitiklį:
\[
n^2+3n+1=(n+1)(n+2)-1.
\]
Todėl
\[
\frac{n^2+3n+1}{n+1} = n+2-\frac{1}{n+1}.
\]
Kadangi $n+2$ yra sveikasis skaičius, visa išraiška bus sveikasis
skaičius tada ir tik tada, kai
\[
\frac{1}{n+1}
\]
yra sveikasis skaičius. Tai įmanoma tik tada, kai
\[
n+1=1 \qquad\text{arba}\qquad n+1=-1.
\]
Taigi
\[
n=0 \qquad\text{arba}\qquad n=-2.
\]

Būtinai patikriname gautus kandidatus pradinėje sąlygoje:
\begin{itemize}
    \item Kai $n=0$: $\frac{0^2+3\cdot 0+1}{0+1} = \frac{1}{1} = 1$ (sveikasis skaičius, tinka).
    \item Kai $n=-2$: $\frac{(-2)^2+3(-2)+1}{-2+1} = \frac{4-6+1}{-1} = \frac{-1}{-1} = 1$ (sveikasis skaičius, tinka).
\end{itemize}

\textbf{Atsakymas:} $\boxed{n=-2,\ 0}$. \hfill $\square$
\end{solution}

\section{Pertvarkykite sąlygą}

Olimpiadiniame uždavinyje informacija dažnai pateikta ne ta forma, kuri
patogiausia sprendimui. Todėl verta klausti:

\[
\text{„Ar galiu tą pačią sąlygą užrašyti naudingesniu būdu?“}
\]
Dažnai padeda:

\begin{itemize}
    \item išskaidymas dauginamaisiais;
    \item bendro daugiklio iškėlimas;
    \item skirtumo kvadratų formulė;
    \item tinkamo skaičiaus pridėjimas arba atėmimas;
    \item visų narių perkėlimas į vieną pusę.
\end{itemize}

\begin{example}
Raskite visas sveikųjų skaičių poras $(x,y)$, tenkinančias
\[
x+y=xy.
\]
\end{example}

\begin{solution}
\textbf{Sprendimas}
Perkeliame narius:
\[
xy-x-y=0.
\]
Pridedame $1$:
\[
xy-x-y+1=1.
\]
Gauname
\[
(x-1)(y-1)=1.
\]
Kadangi $x$ ir $y$ yra sveikieji, turime tik du atvejus:
\[
x-1=1,\qquad y-1=1,
\]
arba
\[
x-1=-1,\qquad y-1=-1.
\]
Todėl
\[
(x,y)=(2,2) \qquad\text{arba}\qquad (x,y)=(0,0).
\]

Patikriname: 
$2+2=2\cdot 2$ (tinka), $0+0=0\cdot 0$ (tinka).

\textbf{Atsakymas:} $\boxed{(0;0) \text{ ir } (2;2)}$. \hfill $\square$
\end{solution}

\section{Trys dažni spąstai}

\subsection{„Patikrinau daug pavyzdžių“}
Tai gali būti puikus būdas \textbf{atspėti} atsakymą, bet neįrodo, kad
jis visada teisingas.

\subsection{„Akivaizdu, kad tai geriausias pasirinkimas“}
Maksimumo arba minimumo uždavinyje neužtenka parodyti vieną gerą
konstrukciją. Reikia dar įrodyti, kad geresnė neįmanoma.

\subsection{„Uždavinys to klausia, vadinasi taip ir yra“}
Uždavinio formuluotė nėra papildoma sąlyga. Jeigu uždavinys prašo rasti
tam tikrą santykį ar kampą, negalima tiesiog manyti, kad jis nekinta.
Tai reikia įrodyti.

\section{Ką daryti, kai nežinote, nuo ko pradėti?}

Užstrigti olimpiadiniame uždavinyje yra normalu. Tikslas nėra visada iš
karto pamatyti sprendimą. Tikslas --- mokėti ieškoti.

Pabandykite eiti tokia tvarka:

\begin{enumerate}
    \renewcommand{\labelenumi}{\textbf{\theenumi.}}

    \item \textbf{Tiksliai suformuluokite tikslą.}
    Ką turite įrodyti arba rasti?

    \item \textbf{Išrašykite, ką žinote.}
    Kokias išvadas iš karto duoda sąlyga?

    \item \textbf{Patikrinkite mažus atvejus.}
    Jie gali padėti pastebėti dėsningumą, bet nepamirškite, kad tai dar
    nėra įrodymas.

    \item \textbf{Ieškokite struktūros.}
    Lyginumas, dalumas, faktorizavimas, simetrija, vienodi reiškiniai,
    keli galimi atvejai.

    \item \textbf{Pabandykite sąlygą perrašyti.}
    Kartais viena tinkama transformacija atskleidžia visą uždavinį.

    \item \textbf{Pagalvokite nuo galo.}
    Ko pakaktų, kad galėtumėte padaryti norimą išvadą?

    \item \textbf{Jeigu vis dar nežinote, užrašykite bent vieną naują
    teisingą išvadą.}
    Nauja išvada dažnai tampa kitu žingsniu.
\end{enumerate}

\section{Kaip turi atrodyti geras olimpiadinis sprendimas?}

Geras sprendimas turi būti ne tik teisingas, bet ir logiškai pilnas.

Prieš baigdami uždavinį patikrinkite:

\begin{itemize}
    \item Ar aišku, ką reiškia visi mano kintamieji?
    \item Ar įrodžiau būtent tą teiginį, kurio buvo prašoma?
    \item Ar kiekviena svarbi išvada yra pagrįsta?
    \item Ar nepakeičiau įrodymo kelių pavyzdžių patikrinimu?
    \item Jei reikėjo rasti visus sprendinius, ar įrodžiau, kad kitų nėra?
    \item Jei ieškojau maksimumo ar minimumo, ar įrodžiau optimalumą?
    \item Ar patikrinau rastus kandidatus pradinėje sąlygoje?
\end{itemize}

Olimpiadoje vertinama ne tai, ką sprendėjas turėjo galvoje, o tai, ką jis
sugebėjo pagrįsti ir užrašyti.

\section{Uždaviniai savarankiškam sprendimui}

Pirmuose uždaviniuose svarbu ne tik gauti atsakymą, bet ir pastebėti,
\textbf{ko trūksta argumentui}. Vėliau pereisime prie savarankiškų
įrodymų.

\subsection{A lygis. Ar tai tikrai įrodymas?}

\begin{enumerate}
    \renewcommand{\labelenumi}{\textbf{\theenumi.}}

    \item Mokinys nori įrodyti, kad $n(n+1)$ yra lyginis kiekvienam
    sveikajam $n$. Jis patikrina $n=1,2,3,4,5$ ir visais atvejais gauna
    lyginį skaičių.
    Ar tai įrodymas? Jei ne, parašykite pilną įrodymą.

    \item Reikia įrodyti teiginį:
    \[
    6\mid n\Longrightarrow 3\mid n.
    \]
    Mokinys užrašo: „Jeigu $3\mid n$, tai $n=3k$. Todėl $n$ dalijasi iš
    $3$.“
    Paaiškinkite, kodėl šis argumentas neįrodo prašomo teiginio, ir
    parašykite teisingą įrodymą.

    \item Mokinys teigia:
    \[
    \text{Jeigu }a^2\text{ dalijasi iš }4,\text{ tai }a\text{ dalijasi iš }4.
    \]
    Raskite kontrpavyzdį. Tada suformuluokite panašų teisingą teiginį.

    \item Uždavinyje reikia rasti visus sveikuosius $x$, kuriems
    \[
    x^2=9.
    \]
    Mokinys parašo: „Tinka $x=3$ ir $x=-3$.“
    Ko trūksta, kad sprendimas būtų pilnas?

    \item Mokinys nori rasti didžiausią sandaugą $xy$, kai $x$ ir $y$
    yra teigiami realieji skaičiai ir
    \[
    x+y=10.
    \]
    Jis parašo: „Akivaizdu, kad geriausia imti $x=y=5$, todėl maksimumas
    yra $25$.“
    Kokia svarbiausia loginė spraga šiame argumente?
\end{enumerate}

\subsection{B lygis. Pirmieji pilni įrodymai}

\begin{enumerate}
    \renewcommand{\labelenumi}{\textbf{\theenumi.}}
    \setcounter{enumi}{5}

    \item Įrodykite, kad lyginio ir nelyginio sveikųjų skaičių suma yra
    nelyginė.

    \item Įrodykite, kad dviejų nelyginių sveikųjų skaičių sandauga yra
    nelyginė.

    \item Įrodykite: jeigu $a+b$ yra nelyginis, tai vienas iš skaičių
    $a,b$ yra lyginis, o kitas --- nelyginis.

    \item Įrodykite: jeigu $n^3$ yra lyginis, tai $n$ yra lyginis.

    \item Įrodykite, kad trijų iš eilės einančių sveikųjų skaičių suma
    dalijasi iš $3$.

    \item Įrodykite, kad tarp bet kurių trijų sveikųjų skaičių yra bent
    du vienodo lyginumo.
\end{enumerate}

\subsection{C lygis. Raskite visus ir pertvarkykite}

\begin{enumerate}
    \renewcommand{\labelenumi}{\textbf{\theenumi.}}
    \setcounter{enumi}{11}

    \item Raskite visus sveikuosius skaičius $n$, kuriems
    \[
    n-2\mid n^2-5n+8.
    \]

    \item Raskite visus sveikuosius skaičius $n$, kuriems
    \[
    n+2\mid n^2+n+1.
    \]

    \item Raskite visas teigiamųjų sveikųjų skaičių poras $(x,y)$,
    tenkinančias
    \[
    xy=x+y+2.
    \]

    \item Raskite visas teigiamųjų sveikųjų skaičių poras $(x,y)$,
    tenkinančias
    \[
    \frac1x+\frac1y=\frac13.
    \]

    \item Raskite visas sveikųjų skaičių poras $(x,y)$, tenkinančias
    \[
    x^2-y^2=15.
    \]
    Įrodykite, kad radote visas poras.
\end{enumerate}

\subsection{D lygis. Pirmasis olimpiadinis iššūkis}

\begin{enumerate}
    \renewcommand{\labelenumi}{\textbf{\theenumi.}}
    \setcounter{enumi}{16}

    \item Įrodykite, kad keturių iš eilės einančių sveikųjų skaičių
    sandauga dalijasi iš $24$.

    \item Tegu $a$ ir $b$ yra sveikieji skaičiai. Įrodykite, kad jeigu
    \[
    a+b
    \]
    ir
    \[
    ab
    \]
    yra lyginiai, tai $a$ ir $b$ abu yra lyginiai.

    \item Įrodykite, kad lygtis
    \[
    x^2-y^2=2
    \]
    neturi sveikųjų sprendinių. 
    \textit{(Užuomina: išskaidykite kairę pusę dauginamaisiais ir pagalvokite apie $(x-y)$ bei $(x+y)$ lyginumą. Kas nutinka, kai lygtį sprendžiame prieštaros būdu?)}

    \item Tarp bet kurių penkių sveikųjų skaičių įrodykite, kad yra du
    tokie, kurių skirtumas dalijasi iš $4$.
\end{enumerate}

\section{Po treniruotės}

Atsakykite į klausimus:

\begin{enumerate}
    \renewcommand{\labelenumi}{\textbf{\theenumi.}}

    \item Kuris uždavinys buvo sunkiausias?
    \item Kuriame uždavinyje ilgiausiai nežinojote, nuo ko pradėti?
    \item Kuriame uždavinyje pavyzdžiai padėjo atspėti idėją?
    \item Kuriame uždavinyje reikėjo įrodyti, kad kitų galimybių nėra?
    \item Kurį uždavinį po savaitės verta išspręsti dar kartą be pagalbos?
\end{enumerate}

\bigskip

\begin{center}
\textbf{Olimpiadininko taisyklė}
\end{center}

Jeigu nežinote viso sprendimo, tai dar nereiškia, kad nieko negalite
padaryti.

Pabandykite:

\[
\boxed{\text{surasti bent vieną naują teisingą išvadą iš to, kas duota.}}
\]

Būtent iš tokių mažų žingsnių dažniausiai ir gimsta visas sprendimas.

\end{document}